\section*{Erd\H{o}s Problem 697: a threshold for congruent divisors}

\paragraph{Statement and status.}
Let \(\delta(m,\alpha)\) be the density of integers divisible by
some \(d\equiv1\pmod m\) with \(1<d<\exp(m^\alpha)\).
The question asks for a constant \(\beta>1\) separating limits
zero and one as \(m\longrightarrow\infty\).
Hall's published resolution gives \(\beta=1/\log2\), as recorded
at \url{https://www.erdosproblems.com/697};
the primary reference is R. R. Hall,
\emph{On some conjectures of Erd\H{o}s in Ast\'erisque, I},
\url{https://doi.org/10.1016/0022-314X(92)90096-8}.
This attempt proves the elementary subcritical range,
an exact finite formula, and a fixed-modulus limiting result.
The sharp threshold argument remains unresolved in this
reconstruction; the known value is attributed rather than
presented as a new proof.

\paragraph{Existence and exact computation of the density.}
For fixed \(m,\alpha\), the admissible divisor set
\[
 D=\{1+jm:j\geq1,\ 1+jm<\exp(m^\alpha)\}
\]
is finite. If it is empty the density is zero.
Otherwise divisibility by a member of \(D\) is periodic modulo
\(L=\operatorname{lcm}(D)\), and
\[
 \delta(m,\alpha)=
 \sum_{\varnothing\ne S\subseteq D}
       \frac{(-1)^{|S|+1}}{\operatorname{lcm}(S)}.               \tag{1}
\]
This follows from inclusion--exclusion and the density
\(1/\operatorname{lcm}(S)\) of common multiples.
In particular there is no question about the existence of
the density before taking the outer limit in \(m\).

\paragraph{A rigorous subcritical bound.}
Let \(J=|D|\). If \(J\geq1\), the union bound gives
\[
 \delta(m,\alpha)\leq\sum_{j=1}^J\frac1{1+jm}
 <\frac1m\sum_{j=1}^J\frac1j
 \leq\frac{1+\log J}{m}
 <\frac{1+m^\alpha}{m}.                                       \tag{2}
\]
The same final bound holds when \(J=0\).
For every fixed \(\alpha<1\), (2) tends to zero.
This proves a genuine part of the requested phase transition,
but leaves the interval \(1\leq\alpha<1/\log2\) untreated.

\paragraph{What can be proved when the modulus is fixed.}
For each fixed \(m\geq2\), the sum of reciprocals of primes
\(q\equiv1\pmod m\) diverges. This is a classical consequence
of the prime number theorem in the fixed progression:
partial summation gives
\(\sum_{q\leq z,\ q\equiv1\ (m)}1/q
  =(\log\log z)/\varphi(m)+O_m(1)\).
Using only such primes as admissible divisors and applying
the Chinese remainder theorem yields
\[
 \delta(m,\alpha)\geq
 1-\prod_{\substack{q<\exp(m^\alpha)\\q\ {\rm prime}\\q\equiv1\ (m)}}
       \left(1-\frac1q\right).
\]
Therefore, for each fixed \(m\geq2\),
\[
 \lim_{\alpha\to\infty}\delta(m,\alpha)=1.                       \tag{3}
\]
The proof uses the fixed-progression prime number theorem
as an explicit external input.
It does not give a bound uniform in \(m\), and so does not
justify exchanging the two limits in (3).

\paragraph{Finite upper and lower certificates.}
For completeness, let
\(\mu=\sum_{d\in D}1/d\) and
\[
 V=\mu+2\sum_{\substack{d,e\in D\\d<e}}\frac1{\operatorname{lcm}(d,e)}.
\]
On residues modulo \(L\), put \(X=\sum_{d\in D}{\bf1}_{d\mid n}\).
Then \(\mathbb EX=\mu\), \(\mathbb EX^2=V\), and Cauchy--Schwarz
applied to \(X{\bf1}_{X>0}\) gives
\[
 \frac{\mu^2}{V}\leq\delta(m,\alpha)\leq\min(1,\mu)
 \quad(D\ne\varnothing).                                      \tag{4}
\]
Both (1) and (4) are exact finite certificates that do not
assume independence of composite-divisor events.

\paragraph{Remaining gap.}
The uniform distribution of these highly dependent composite
divisors is the missing step. Neither (2), (3), nor the
second-moment bound (4) proves Hall's threshold or the limit
one for any fixed \(\alpha>1/\log2\) as \(m\) grows.
This attempt remains unresolved beyond the partial results stated.

\paragraph{Completion estimate.}
\textbf{COMPLETION: 40\%}
